\section{The literal projective construction}
Theorem 6.1 of Leemans--Stokes--Tranchida takes the old objects of
$P=\operatorname{PG}(n-1,K)$ and, for each $\lambda\in K\setminus F$,
adds the ordered collinear quadruples with cross-ratio $\lambda$ as a
new type $Q_\lambda$. An old point is incident with a quadruple exactly
when it is one of its four entries; an old higher-dimensional subspace
is incident exactly when it contains the four points. Any two new
objects of different types are incident. The old shadow therefore
remembers the unordered support, not its ordering.

\begin{theorem}[Exact obstruction over four elements]\label{A126.Projective}
For the literal construction with $n=3$, $F=\mathbb F_2$ and
$K=\mathbb F_4$, the pointwise-old-object kernel in the type-preserving
group is $\Sym(12)^{210}$. Hence the literal system is not a
representation of $\operatorname{PGL}(3,4)$. The same instance refutes
the literal root-type truncation of Corollary 6.2.
\end{theorem}
\begin{proof}
Write $K=\mathbb F_2[a]/(a^2+a+1)$. The old plane has 21 points and
21 lines, each line containing five points. For any four-subset $A$ of
a line $\ell$, all 24 orderings have old shadow $A\cup\{\ell\}$.
Their cross-ratios are neither zero nor one, so they belong to
$\{a,a+1\}$. Interchanging the last two positions inverts the
cross-ratio. Since $a^{-1}=a+1$, this is a bijection between the two
sets of orderings. Each type therefore has exactly twelve objects
with this shadow.

There are $21\binom54=105$ distinct such shadows, at each of two
new types. the complete shadow-kernel theorem of the common paper gives exactly
$\Sym(12)^{210}$ as the pointwise kernel. In particular the
type-preserving group contains $\Sym(12)$, of order 479001600,
whereas
\[
 |\operatorname{PGL}(3,4)|=
 \frac{(4^3-1)(4^3-4)(4^3-4^2)}{4-1}=60480.
\]
Thus even the abstract group orders preclude the proposed isomorphism.
There are $42+2\cdot105\cdot12=2562$ objects in the literal model.

For Corollary 6.2 choose $\alpha=a$. Its minimal polynomial over
$\mathbb F_2$ is $X^2+X+1$, with precisely $a,a+1$ as roots in $K$.
For $n=3$ the displayed truncation keeps both old types and both new
types. It is exactly the same counterexample.
\end{proof}

The source proof's uniqueness assertion is the failed step. The always
valid double interchange
$(p_1,p_2,p_3,p_4)\mapsto(p_2,p_1,p_4,p_3)$ already gives different
objects with the same cross-ratio and old shadow. The finite example
above supplies the complete kernel and an order contradiction; it does
not depend on identifying the full correlation group.

\begin{proposition}[Collapsing twins does not restore the field label]
In this four-element example the image of restriction is all $K(P)$.
The nontrivial field automorphism has a type-preserving lift to the
literal extension.
\end{proposition}
\begin{proof}
At both new types, every four-subset on every line has the same
multiplicity twelve. Every old type-preserving automorphism preserves
those multiplicities, so $H=K(P)$ in that theorem. Explicitly,
apply Frobenius $x\mapsto x^2$ to all coordinates and all quadruple
entries, then interchange the last two entries of each quadruple.
The new cross-ratio is $(\lambda^2)^{-1}=\lambda$ for
$\lambda\in\{a,a+1\}$. The two operations commute and each squares
to the identity; their composite is an involutory type-preserving
lift of Frobenius.
\end{proof}

The ordered-position proposition in the common paper gives a role-preserving ordered-tuple
encoding. By itself it does not give a natural projective representation:
recognizability under all correlations, extension of duality and the palette
action are additional requirements. The distinct-arm construction below
discharges them over finite fields.
For finite groups $L$, including each finite projective group, the
complete finite normal-pair theorem of D\'iaz Ramos--Molinier--Viruel,
proved in the first application, supplies a representation by substituting
$(G,H)=(\Aut(L),\operatorname{Inn}(L))$. This is that theorem's
existence consequence, not a verification of the literal projective model.

Two further scope points are visible in the source definitions. Its
preliminary definition requires a finite palette, whereas $K\setminus F$
can be infinite in Section 6. Also the full semilinear group uses all
field automorphisms; identifying it with $\operatorname{Gal}(K/F)$
requires those automorphisms to fix $F$. The finite $\mathbb F_4/\mathbb F_2$
counterexample satisfies both favorable conditions, so neither
qualification removes the obstruction. The orthogonal, degree-two and
Suzuki variants in Section 6.2 are proposed investigations, explicitly
without the necessary construction details in the source; no natural
model for those suggestions is claimed here.

\begin{proposition}[A finite-palette obstruction over a rational function field]
\label{A126.InfiniteField}
For $n\ge3$, the group $L=\operatorname{PGL}(n,\mathbb Q(t))$ has no
exact incidence representation with finitely many types, even if infinitely
many objects are allowed.
\end{proposition}
\begin{proof}
For each integer $m$, substitution $t\mapsto t+m$ is a field automorphism
and therefore acts entrywise as an automorphism $\psi_m$ of $L$.
These automorphisms have distinct images modulo $\operatorname{Inn}(L)$.
Indeed $\psi_m\psi_k^{-1}=\psi_{m-k}$. If $\psi_h$ were inner, the
projective classes of $D=\operatorname{diag}(t,1,\ldots,1)$ and
$\operatorname{diag}(t+h,1,\ldots,1)$ would be conjugate. Thus one
matrix $\operatorname{diag}(t+h,1,\ldots,1)$ would be conjugate in
$\operatorname{GL}(n,\mathbb Q(t))$ to a scalar multiple $uD$. The repeated eigenvalue has multiplicity
$n-1\ge2$ and the other eigenvalue has multiplicity one, so it forces
$u=1$. The remaining eigenvalue forces $h=0$. Thus
$\operatorname{Out}(L)$ is infinite.

For any incidence system with a finite palette $I$, the homomorphism
from correlations to $\Sym(I)$ has kernel precisely the type-preserving
group. Consequently $C/K$ is finite. An exact pair isomorphism with
$(\Aut(L),\operatorname{Inn}(L))$ would identify this finite quotient
with $\operatorname{Out}(L)$, a contradiction.
\end{proof}
This concerns the definition in the source, whose palette is finite.
It is not an obstruction to a different theory allowing infinitely
many types, nor does it rule out weak representations.

\subsection{A faithful finite-field construction}
Here is a positive replacement of the intended finite-field mechanism.
It changes the attachment and includes both point and dual pencils, so
it is distinct from the literal source construction just refuted.

Let $q=p^f$ be a prime power and $n\ge3$. Take the full proper nonzero
subspace incidence system $P=\operatorname{PG}(n-1,q)$, with dimensions
as types. If $f>1$, choose either
\[
R=\mathbb F_q\setminus\mathbb F_p,
\quad\text{or}\quad
R=\{\alpha,\alpha^p,\ldots,\alpha^{p^{f-1}}\}
\quad\text{where }\mathbb F_p(\alpha)=\mathbb F_q.
\]
For $f=1$ put $R=\varnothing$. For every $\lambda\in R$ take all ordered
quadruples of distinct collinear points with cross-ratio $\lambda$,
and all ordered quadruples of hyperplanes in a pencil with cross-ratio
$\lambda$. The latter is the usual cross-ratio on the corresponding
projective line in the dual vector space. The center color is
$\lambda$ in both cases. Apply the rigid-arms theorem of the common paper with the arm lengths
$2,3,4,5$. Write $Y_{n,q,R}$ for the resulting incidence system.

\begin{theorem}[Finite projective representation with ordered attachments]
\label{A126.OrderedProjective}
For $n\ge3$, restriction identifies the pair of the system above with
\[
(C(Y_{n,q,R}),K(Y_{n,q,R}))
 \cong
(\operatorname{P\Gamma L}(n,q)\rtimes\langle\delta\rangle,
 \operatorname{PGL}(n,q)),
\]
where $\delta$ is the standard point--hyperplane duality of order two.
With the retained classical identification of the full projective-group
automorphism group, this is an exact representation of
$\operatorname{PGL}(n,q)$. A finite incidence geometry with the same pair
then follows from the complete cleaning and chamber theorems of the
common paper.
\end{theorem}
\begin{proof}
If $f=1$, no objects are added, and the assertion is the classical
fundamental theorem for finite projective spaces, including duality.
Suppose $f>1$, so $q\ge4$. Every old dimension type has size
\[
 {n\brack d}_q=
 \prod_{a=0}^{d-1}\frac{q^{n-a}-1}{q^{d-a}-1}
 \equiv1\pmod p.
\]
The denominators are nonzero modulo $p$, which justifies the displayed
congruence of integers. There are $L={n\brack2}_q$ point lines and
also $L$ hyperplane pencils. For a fixed cross-ratio other than zero
and one, a projective line has exactly $(q+1)q(q-1)$ ordered quadruples:
choose the first three distinct points, and the cross-ratio determines
the fourth uniquely. Thus each center color has
\[
 M=2L(q+1)q(q-1)
\]
objects; each internal arm color has $M|R|$ objects. Every new type
size is divisible by $p$. Since a correlation permutes entire type
classes and preserves their sizes, it preserves precisely the union
of old types, characterized by size not divisible by $p$.
The recognition hypothesis of the rigid-arms theorem is proved.

We retain the classical fundamental theorem of projective geometry:
the correlations of the full subspace incidence system are the
semilinear projectivities and their composites with duality. A
semilinear transformation with field part $\sigma$ sends cross-ratios
to $\sigma(\lambda)$. Both choices of $R$ are invariant under every
field automorphism. The duality sends an ordered point quadruple to
an ordered hyperplane-pencil quadruple with the same cross-ratio,
and conversely. This can be seen in coordinates by representing a
pencil by its two-dimensional space of linear forms; the standard
duality uses the same coefficient vectors. Hence every old correlation
preserves the tuple family and permutes its center colors. The rigid-arm
theorem supplies its unique extension and rules out further ones.

The internal colors are fixed by every extension. Fixing the old
colors rules out duality because $n\ge3$ and points and hyperplanes
are different types. The remaining action has a field part $\sigma$;
fixing every center color means fixing $R$ pointwise. If $R$ is the
orbit of $\alpha$, this fixes the field generator $\alpha$, hence
$\sigma=1$. If $R=\mathbb F_q\setminus\mathbb F_p$, fixing it
pointwise also fixes the prime field and therefore all of
$\mathbb F_q$. Consequently the palette kernel is exactly
$\operatorname{PGL}(n,q)$. The converse inclusion follows because
projectivities preserve cross-ratios and dimensions.

Finally we retain the standard finite classical-group identification
$\Aut(\operatorname{PGL}(n,q))=
\operatorname{P\Gamma L}(n,q)\rtimes\langle\delta\rangle$ and
$\operatorname{Inn}(\operatorname{PGL}(n,q))=
\operatorname{PGL}(n,q)$ for $n\ge3$. These are the group-theoretic
inputs identifying the normal pair just proved with the desired pair.
The two common conversion theorems apply to the finite properly colored
graph $Y_{n,q,R}$ and preserve both groups, giving the geometric version.
\end{proof}

The prime-field generation condition in the smaller label set is
intentional. Generating $K$ over a larger prescribed subfield need not
force all automorphisms of $K$ fixing the chosen labels to be trivial.
The theorem asserts the displayed finite-field construction only. It
makes no arbitrary infinite-field representation claim, and it does
not identify its new objects with the source's old quadruple objects.
